\documentclass[a4paper,11pt]{article}
\usepackage[utf8]{inputenc}
\usepackage[T1]{fontenc}
\usepackage[brazilian]{babel}
\usepackage[margin=2.5cm]{geometry}
\usepackage{hyperref}
\usepackage{longtable}
\usepackage{booktabs}
\usepackage{array}
\usepackage{xcolor}

\hypersetup{
  colorlinks=true,
  linkcolor=blue!70!black,
  urlcolor=blue!70!black,
  citecolor=blue!70!black
}

\newcommand{\doi}[1]{\href{https://doi.org/#1}{\texttt{#1}}}

\title{\textbf{Referências-chave por Variável do Modelo Hexadimensional (V1--V6)}\\[6pt]
\large Índice de Salvaguarda Biocultural (ISB)}
\author{Diego Vidal --- Compilação para leitura dirigida}
\date{Março 2026}

\begin{document}
\maketitle
\tableofcontents
\newpage

%% ====================================================================
\section{V1 --- Erosão Intergeracional do Conhecimento Tradicional}
%% ====================================================================

\begin{longtable}{>{\raggedright\arraybackslash}p{0.75\textwidth} c}
\toprule
\textbf{Referência} & \textbf{Cit.} \\
\midrule
\endfirsthead
\toprule
\textbf{Referência} & \textbf{Cit.} \\
\midrule
\endhead

Reyes-García, V. et al. (2013). Evidence of traditional knowledge loss among a contemporary indigenous society. \textit{Evolution and Human Behavior}, 34(4), 249--257.\newline
\doi{10.1016/j.evolhumbehav.2013.01.002} & 349 \\
\addlinespace

Gómez-Baggethun, E. \& Reyes-García, V. (2013). Reinterpreting change in traditional ecological knowledge. \textit{Human Ecology}, 41, 643--647.\newline
\doi{10.1007/s10745-013-9577-9} & 300 \\
\addlinespace

Tang, R. \& Gavin, M.C. (2016). A classification of threats to traditional ecological knowledge and conservation responses. \textit{Conservation and Society}, 14(1), 57--70.\newline
\doi{10.4103/0972-4923.182799} & 169 \\
\addlinespace

Reyes-García, V. et al. (2014). Resilience of traditional knowledge systems: the case of agricultural knowledge in home gardens of the Iberian Peninsula. \textit{Global Environmental Change}, 24, 223--231.\newline
\doi{10.1016/j.gloenvcha.2013.11.004} & 176 \\
\addlinespace

Reyes-García, V. et al. (2006). Cultural, practical, and economic value of wild plants: a quantitative study in the Bolivian Amazon. \textit{Journal of Ethnobiology and Ethnomedicine}, 2, 21.\newline
\doi{10.1186/1746-4269-2-21} & 127 \\
\addlinespace

Hedges, S. \& Kipila, C. (2020). Intergenerational erosion of traditional ecological knowledge and its implications for Maasai rangeland management. \textit{Journal of Ethnobiology}, 40(4), 535--551.\newline
\doi{10.2993/0278-0771-40.4.535} & 27 \\
\addlinespace

Gruberg, H. et al. (2022). Traditional ecological knowledge erosion among the Tsimane'. \textit{Human Ecology}, 50, 1053--1066.\newline
\doi{10.1007/s10745-022-00375-9} & 23 \\
\addlinespace

Zent, S. \& Maffi, L. (2009). Final report on indicator No.~2: Vitality Index of Traditional Environmental Knowledge (VITEK). GEO-4 / IUCN.\newline
\href{https://www.terralingua.org/vitek/}{terralingua.org/vitek} & 46 \\

\bottomrule
\end{longtable}

%% ====================================================================
\section{V2 --- Complexidade Biocultural}
%% ====================================================================

\begin{longtable}{>{\raggedright\arraybackslash}p{0.75\textwidth} c}
\toprule
\textbf{Referência} & \textbf{Cit.} \\
\midrule
\endfirsthead
\toprule
\textbf{Referência} & \textbf{Cit.} \\
\midrule
\endhead

Maffi, L. (2001). \textit{On Biocultural Diversity: Linking Language, Knowledge, and the Environment}. Washington: Smithsonian Institution Press. ISBN 1-56098-905-X.\newline
\href{https://www.google.com/search?q=Maffi+2001+On+Biocultural+Diversity+Smithsonian}{Buscar} & 700+ \\
\addlinespace

Pretty, J. et al. (2009). The intersections of biological diversity and cultural diversity: towards integration. \textit{Conservation and Society}, 7(2), 100--112.\newline
\doi{10.4103/0972-4923.58642} & 400+ \\
\addlinespace

Whitney, C. et al. (2018). Ethnobotany and agrobiodiversity: valuation of plants in the homestead gardens of southwestern Uganda. \textit{Ethnobiology Letters}, 9(2), 164--178.\newline
\href{https://www.jstor.org/stable/26551443}{JSTOR} & 39 \\
\addlinespace

Gepts, P. (2023). Biocultural diversity and crop improvement. \textit{Current Opinion in Plant Biology}, 71, 102299.\newline
\doi{10.1016/j.pbi.2022.102299} & 33 \\
\addlinespace

Fernández-Llamazares, Á. et al. (2021). Scientists' warning to humanity on threats to Indigenous and local knowledge systems. \textit{Journal of Ethnobiology}, 41(2), 144--169.\newline
\doi{10.2993/0278-0771-41.2.144} & 350 \\
\addlinespace

Fernández-Llamazares, Á. et al. (2020). The role of Indigenous Peoples and local communities in ecological restoration. \textit{Restoration Ecology}, 29(2), e13174.\newline
\doi{10.1111/rec.13174} & $\sim$100 \\

\bottomrule
\end{longtable}

%% ====================================================================
\section{V3 --- Singularidade Territorial}
%% ====================================================================

\begin{longtable}{>{\raggedright\arraybackslash}p{0.75\textwidth} c}
\toprule
\textbf{Referência} & \textbf{Cit.} \\
\midrule
\endfirsthead
\toprule
\textbf{Referência} & \textbf{Cit.} \\
\midrule
\endhead

Fernández-Llamazares, Á. et al. (2021). Scientists' warning to humanity on threats to ILK systems. \textit{Journal of Ethnobiology}, 41(2), 144--169.\newline
\doi{10.2993/0278-0771-41.2.144} & 350 \\
\addlinespace

Zimmerer, K.S. et al. (2018). Mountain ecology, remoteness, and the rise of agrobiodiversity: tracing the geographic spaces of human-environment knowledge. \textit{Annals of the AAG}, 108(2), 355--375.\newline
\doi{10.1080/24694452.2017.1354069} & $\sim$80 \\
\addlinespace

Rayne, A. et al. (2022). Integrating place-based knowledge for culturally significant species recovery. \textit{Evolutionary Applications}, 15(4), 630--644.\newline
\doi{10.1111/eva.13367} & $\sim$30 \\
\addlinespace

Cámara-Leret, R. \& Bascompte, J. (2021). Language extinction triggers the loss of unique medicinal knowledge. \textit{PNAS}, 118(24), e2103683118.\newline
\doi{10.1073/pnas.2103683118} & 200+ \\
\addlinespace

Berkes, F. (2012). \textit{Sacred Ecology: Traditional Ecological Knowledge and Resource Management}. 3rd ed. Routledge.\newline
\doi{10.4324/9780203123843} & 5000+ \\

\bottomrule
\end{longtable}

%% ====================================================================
\section{V4 --- Status de Documentação}
%% ====================================================================

\begin{longtable}{>{\raggedright\arraybackslash}p{0.75\textwidth} c}
\toprule
\textbf{Referência} & \textbf{Cit.} \\
\midrule
\endfirsthead
\toprule
\textbf{Referência} & \textbf{Cit.} \\
\midrule
\endhead

Maina, C.K. (2012). Traditional knowledge management and preservation: intersections with Library and Information Science. \textit{International Information and Library Review}, 44(1), 13--27.\newline
\doi{10.1016/j.iilr.2012.01.001} & 119 \\
\addlinespace

Poorna, R.L., Mymoon, M. \& Hariharan, A. (2014). Preservation and protection of traditional knowledge --- diverse documentation initiatives across the globe. \textit{Current Science}, 107(8), 1240--1246.\newline
\href{https://www.jstor.org/stable/24103468}{JSTOR / Current Science 107(8)} & 88 \\
\addlinespace

Balogun, T. \& Kalusopa, T. (2021). Digital preservation of indigenous knowledge systems: a framework for South Africa. \textit{Records Management Journal}, 31(3), 304--321.\newline
\doi{10.1108/RMJ-12-2020-0042} & 54 \\
\addlinespace

Oruç, S. (2022). Documenting Indigenous oral traditions: copyright and cultural heritage considerations. \textit{Journal of Copyright in Education and Librarianship}, 6(1).\newline
\href{https://doi.org/10.17161/jcel.v6i1}{JCEL 6(1)} & 14 \\
\addlinespace

Nakata, M. \& Langton, M. (2005). Australian indigenous knowledge and libraries. \textit{Australian Academic \& Research Libraries}, 36(2), 2--14.\newline
\doi{10.1080/00048623.2005.10755292} & 100+ \\

\bottomrule
\end{longtable}

%% ====================================================================
\section{V5 --- Vulnerabilidade Jurídica}
%% ====================================================================

\begin{longtable}{>{\raggedright\arraybackslash}p{0.75\textwidth} c}
\toprule
\textbf{Referência} & \textbf{Cit.} \\
\midrule
\endfirsthead
\toprule
\textbf{Referência} & \textbf{Cit.} \\
\midrule
\endhead

Correa, C.M. (2001). \textit{Traditional Knowledge and Intellectual Property}. Geneva: QUNO.\newline
\href{https://quakerservice.ca/wp-content/uploads/2011/07/TKcol3.pdf}{PDF livre (QUNO)} & 223 \\
\addlinespace

Dutfield, G. (2004). \textit{Intellectual Property, Biogenetic Resources and Traditional Knowledge}. London: Earthscan/Routledge.\newline
\doi{10.4324/9781849770859} & 600+ \\
\addlinespace

Cottier, T. \& Panizzon, M. (2004). Legal perspectives on traditional knowledge: the case for intellectual property protection. \textit{J. International Economic Law}, 7(2), 371--399.\newline
\doi{10.1093/jiel/7.2.371} & 176 \\
\addlinespace

Bavikatte, K.S. \& Robinson, D.F. (2011). Towards a people's history of the law: biocultural jurisprudence and the Nagoya Protocol on access and benefit sharing. \textit{Law, Environment and Development Journal}, 7, 35.\newline
\href{http://www.lead-journal.org/content/11035.pdf}{LEAD Journal 7/35 (PDF livre)} & 137 \\
\addlinespace

Osei-Tutu, J.J. (2011). A sui generis regime for traditional knowledge: the cultural divide in intellectual property law. \textit{Marquette Intellectual Property Law Review}, 15(1), 147.\newline
\href{https://scholarship.law.marquette.edu/iplr/vol15/iss1/6/}{Marquette IPLR 15(1)} & 138 \\
\addlinespace

Nemogá, G.R., Appasamy, A. et al. (2022). Protecting indigenous and local knowledge through a biocultural diversity framework. \textit{The Journal of Environment \& Development}, 31(4), 385--413.\newline
\doi{10.1177/10704965221104781} & 58 \\
\addlinespace

Swiderska, K. (2006). \textit{Banishing the biopirates: a new approach to protecting traditional knowledge}. IIED Gatekeeper Series 129.\newline
\href{https://www.jstor.org/stable/resrep01338}{JSTOR / IIED Gatekeeper 129} & 37 \\
\addlinespace

Swiderska, K. (2006). \textit{Protecting Traditional Knowledge: A Framework Based on Customary Laws and Bio-cultural Heritage}. IIED.\newline
\href{https://www.iied.org/sites/default/files/pdfs/migrate/G01069.pdf}{PDF livre (IIED)} & 44 \\
\addlinespace

Robinson, D.F. (2014). \textit{Biodiversity, Access and Benefit-Sharing: Global Case Studies}. Routledge.\newline
\doi{10.4324/9781315882819} & 89 \\
\addlinespace

Dagne, T. (2014). Protecting traditional knowledge in international intellectual property law: imperatives for protection and choice of modalities. \textit{John Marshall Review of Intellectual Property Law}, 14(1), 1--56.\newline
\href{https://repository.law.uic.edu/ripl/vol14/iss1/1/}{JMRIPL 14(1)} & 34 \\

\bottomrule
\end{longtable}

%% ====================================================================
\section{V6 --- Organização Social e Governança}
%% ====================================================================

\begin{longtable}{>{\raggedright\arraybackslash}p{0.75\textwidth} c}
\toprule
\textbf{Referência} & \textbf{Cit.} \\
\midrule
\endfirsthead
\toprule
\textbf{Referência} & \textbf{Cit.} \\
\midrule
\endhead

Berkes, F. (2004). Rethinking community-based conservation. \textit{Conservation Biology}, 18(3), 621--630.\newline
\doi{10.1111/j.1523-1739.2004.00077.x} & 3289 \\
\addlinespace

Olsson, P., Folke, C. \& Berkes, F. (2004). Adaptive comanagement for building resilience in social--ecological systems. \textit{Environmental Management}, 34, 75--90.\newline
\doi{10.1007/s00267-003-0101-7} & 2881 \\
\addlinespace

Berkes, F. (2007). Community-based conservation in a globalized world. \textit{PNAS}, 104(39), 15188--15193.\newline
\doi{10.1073/pnas.0702098104} & 1698 \\
\addlinespace

Berkes, F. \& Turner, N.J. (2006). Knowledge, learning and the evolution of conservation practice for social-ecological system resilience. \textit{Human Ecology}, 34, 479--494.\newline
\doi{10.1007/s10745-006-9008-2} & 878 \\
\addlinespace

Folke, C. (2004). Traditional knowledge in social--ecological systems (editorial). \textit{Ecology and Society}, 9(3), 7.\newline
\href{https://www.ecologyandsociety.org/vol9/iss3/art7/}{Ecology and Society 9(3):7} & 546 \\
\addlinespace

Berkes, F. (2017). Environmental governance for the Anthropocene? Social-ecological systems, resilience, and collaborative learning. \textit{Sustainability}, 9(7), 1232.\newline
\doi{10.3390/su9071232} & 549 \\
\addlinespace

Rathwell, K.J., Armitage, D. \& Berkes, F. (2015). Bridging knowledge systems to enhance governance of the environmental commons: a typology of settings. \textit{International Journal of the Commons}, 9(2), 751--773.\newline
\doi{10.18352/ijc.526} & 201 \\
\addlinespace

Ruiz-Mallén, I. \& Corbera, E. (2013). Community-based conservation and traditional ecological knowledge: implications for social-ecological resilience. \textit{Ecology and Society}, 18(4), 12.\newline
\doi{10.5751/ES-05867-180412} & 313 \\
\addlinespace

Reyes-García, V. \& Benyei, P. (2018). Traditional agricultural knowledge as a commons. In \textit{Routledge Handbook of Food as a Commons}. Routledge.\newline
\doi{10.4324/9781315161495-11} & 36 \\
\addlinespace

Luo, Y. et al. (2024). Village leadership, social networks and collective actions in indigenous communities: Hani rice terraces, China. \textit{Journal of Rural Studies}, 105, 103187.\newline
\doi{10.1016/j.jrurstud.2024.103187} & 35 \\
\addlinespace

Nakamura, N. \& Kanemasu, Y. (2020). Traditional knowledge, social capital, and community response to a disaster: resilience of remote communities in Fiji. \textit{Regional Environmental Change}, 20, 25.\newline
\doi{10.1007/s10113-020-01613-w} & 70 \\

\bottomrule
\end{longtable}

\end{document}
